\documentclass[11pt]{article}

\usepackage[margin=1in]{geometry}
\usepackage{amsmath,amssymb,amsthm,mathtools}
\usepackage{microtype}
\usepackage{enumitem}
\usepackage[hidelinks]{hyperref}
\hypersetup{
  pdftitle={A Counterexample to a Character Count under Coprime Action},
  pdfauthor={Eric Hou}
}

\newtheorem{theorem}{Theorem}[section]
\newtheorem{proposition}[theorem]{Proposition}
\newtheorem{lemma}[theorem]{Lemma}
\newtheorem{corollary}[theorem]{Corollary}
\theoremstyle{definition}
\newtheorem{definition}[theorem]{Definition}
\theoremstyle{remark}
\newtheorem{remark}[theorem]{Remark}

\newcommand{\F}{\mathbf F}
\newcommand{\Irr}{\operatorname{Irr}}
\newcommand{\Ind}{\operatorname{Ind}}
\newcommand{\Res}{\operatorname{Res}}
\newcommand{\supp}{\operatorname{supp}}
\newcommand{\wt}{\operatorname{wt}}
\newcommand{\one}{\mathbf 1}

\title{A Counterexample to a Character Count under Coprime Action}
\author{Eric Hou}
\date{}

\begin{document}
\maketitle

\begin{abstract}
Let a finite group $A$ act coprimely on a finite group $G$, and set
$C=C_G(A)$.  For
\[
N_A(G)=\{\chi\in\Irr(G):\chi^a=\chi\text{ for all }a\in A,
\ \chi(c)\neq 0\text{ for all }c\in C\},
\]
we consider whether $|N_A(G)|=|C/C'|$ always holds.  We construct a
counterexample with $A\cong C_{21}$ and
$G\cong C_2^{128}\rtimes C_2^7$.  Here $C\cong C_2^{12}$, while the
little-groups method identifies the $A$-invariant irreducible characters
with the $4096$ $A$-fixed Boolean functions on
$\F_4^2\oplus\F_8$.  Their values on $C$ reduce to correlations of
six-bit radial functions on $\F_4^2$, and the resulting binary systems
give
\[
|N_A(G)|=1728<4096=|C/C'|.
\]
Three exact computer calculations independently check the count but are
not used in the proof.
\end{abstract}

\section{Introduction}

Let $A$ act by automorphisms on a finite group $G$, with
$(|A|,|G|)=1$, and set $C=C_G(A)$.  We consider whether the number of
$A$-invariant irreducible characters of $G$ that are nonzero on $C$
must equal the number of linear characters of $C$:
\begin{equation}\label{eq:question}
 |N_A(G)|\stackrel{?}{=}|C/C'|,
 \qquad
 N_A(G)=\{\chi\in\Irr^A(G):\chi(c)\neq0\text{ for all }c\in C\}.
\end{equation}
This question, proposed by G. Navarro, appears as Problem 21.100 in the
21st issue of the Kourovka Notebook \cite{kourovka}.

We construct a cyclic group $A$ of order $21$ acting on a $2$-group
$G=B\rtimes V$, where $V\cong C_2^7$ and $B$ is the additive group of
all Boolean functions on $V$.  The action on $V$ has no nonzero fixed
point and has twelve orbits, giving
$C_G(A)=B^A\cong C_2^{12}$.

The little-groups method gives an explicit parameterization of
$\Irr(G)$.  Each $A$-stable translation orbit in $B$ has a unique
$A$-fixed representative, and the only invariant character of its
translation stabilizer is the trivial character.  Thus $\Irr^A(G)$ is
indexed by $B^A$.  Writing an $A$-fixed Boolean function on
$\F_4^2\oplus\F_8$ as a pair of six-bit radial functions on
$\F_4^2$ reduces the character values on $C$ to two binary correlation
functions.  The resulting zero criterion gives $|N_A(G)|=1728$.

The construction is not claimed to have minimal order, and no claim is
made about operator groups subject to additional restrictions.

\section{The construction and its fixed subgroup}

Let
\[
 U=\F_4^2,\qquad W=\F_8,\qquad V=U\oplus W
\]
as additive groups.  Choose $\omega\in\F_4^\times$ of order $3$ and
$\zeta\in\F_8^\times$ of order $7$.  Define
\[
 T(x_1,x_2,y)=(\omega x_1,\omega x_2,\zeta y).
\]
Then $T$ is $\F_2$-linear and has order $21$.  Moreover,
\begin{equation}\label{eq:no-fixed-V}
 C_V(T)=0,
\end{equation}
because $\omega-1$ and $\zeta-1$ are nonzero field elements.

Let
\[
 B=\{f:V\to\F_2\}
\]
with pointwise addition.  For $t\in V$, define
\[
 (\tau_t f)(x)=f(x+t).
\]
The additive group $V$ acts on $B$ through these translations.  Put
\[
 G=B\rtimes V,
 \qquad
 (f,t)(g,s)=(f+\tau_tg,t+s).
\]
Thus $|G|=2^{128}2^7=2^{135}$.

Let $A=\langle a\rangle\cong C_{21}$.  Define
\[
 a\cdot(f,t)=(f\circ T^{-1},Tt).
\]
This is an automorphism of $G$, since
\[
 (\tau_t f)\circ T^{-1}=\tau_{Tt}(f\circ T^{-1}).
\]
Its restriction to $V$ is $T$, so the action is faithful.  Since
$|A|=21$ and $|G|$ is a power of $2$, the action is coprime.

\begin{proposition}\label{prop:fixed-subgroup}
For the action above,
\[
 C_G(A)=B^A\cong C_2^{12}.
\]
In particular,
\[
 |C_G(A)/C_G(A)'|=4096.
\]
\end{proposition}

\begin{proof}
If $(f,t)\in G$ is fixed by $a$, then $Tt=t$.  Equation
\eqref{eq:no-fixed-V} gives $t=0$, so $C_G(A)=B^A$.

It remains to count the $A$-orbits on $V$.  The zero vector is one
orbit.  The nonzero vectors in $U$ split into the five sets
$L_i^\times$, where $L_1,\ldots,L_5$ are the one-dimensional
$\F_4$-subspaces of $U$; each has size $3$ and is one orbit under
multiplication by $\omega$.  The set $\{0\}\times W^\times$ is one
orbit of size $7$.  Finally, each
\[
 L_i^\times\times W^\times
\]
is an orbit of size $21$, since the powers of $(\omega,\zeta)$ run
through $C_3\times C_7$.

The orbit sizes are therefore
\[
 1,3,3,3,3,3,7,21,21,21,21,21.
\]
An element of $B$ is fixed by $A$ exactly when it is constant on each
of these twelve orbits.  Hence $B^A\cong C_2^{12}$.  This group is
abelian, so its derived subgroup is trivial and the final assertion
follows.
\end{proof}

\section{The invariant irreducible characters}

For $u,f\in B$, let
\[
 \langle u,f\rangle=\sum_{x\in V}u(x)f(x)\in\F_2
\]
and define
\[
 \lambda_u(f)=(-1)^{\langle u,f\rangle}.
\]
The map $u\mapsto\lambda_u$ identifies $B$ with $\Irr(B)$.

For $u\in B$, write
\[
 V_u=\{t\in V:\tau_tu=u\},
 \qquad I_u=B\rtimes V_u.
\]
The subgroup $I_u$ is the inertia subgroup of $\lambda_u$ in $G$.
For $\mu\in\Irr(V_u)$, define a linear character of $I_u$ by
\[
 \theta_{u,\mu}(f,t)=\lambda_u(f)\mu(t).
\]
This is well defined because $V_u$ fixes $\lambda_u$.  Put
\[
 \chi_{u,\mu}=\Ind_{I_u}^G\theta_{u,\mu}.
\]

\begin{proposition}\label{prop:all-irr}
Choose one $u$ from each translation orbit of $V$ on $B$.  Then
\[
 \Irr(G)=
 \{\chi_{u,\mu}:u\text{ is a chosen representative},\
 \mu\in\Irr(V_u)\}.
\]
Every character in this set has degree $[V:V_u]$.
\end{proposition}

\begin{proof}
The subgroup $I_u$ is normal in $G$, since it is the inverse image of
$V_u$ under $G\to G/B=V$ and $V$ is abelian.  Mackey's inner-product
formula gives
\[
 \langle\chi_{u,\mu},\chi_{u,\nu}\rangle_G
 =\sum_{gI_u\in G/I_u}
 \langle\theta_{u,\mu},\theta_{u,\nu}^{\,g}\rangle_{I_u}.
\]
If $g\notin I_u$, the restrictions of
$\theta_{u,\mu}$ and $\theta_{u,\nu}^{\,g}$ to $B$ are distinct,
because $g$ sends $\lambda_u$ to a different translate.  Those terms
vanish.  The identity coset contributes $1$ when $\mu=\nu$ and $0$
otherwise.  Thus the characters $\chi_{u,\mu}$ are irreducible and are
distinct for fixed $u$.  Characters arising from different translation
orbits have disjoint sets of constituents on restriction to $B$, so
they are also distinct.

For one translation orbit represented by $u$, there are $|V_u|$
characters, each of degree $[V:V_u]$.  Their contribution to the sum of
squared degrees is
\[
 |V_u|[V:V_u]^2=|V|\,|V\cdot u|.
\]
Summing over all translation orbits in $B$ gives
\[
 \sum_{u,\mu}\chi_{u,\mu}(1)^2
 =|V|\sum_{\mathcal O\subseteq B}|\mathcal O|
 =|V||B|=|G|.
\]
Since these distinct irreducible characters account for the full sum
of squared degrees, they exhaust $\Irr(G)$.
\end{proof}

\begin{lemma}\label{lem:fixed-representative}
Every $A$-stable translation orbit in $B$ contains exactly one element
of $B^A$.
\end{lemma}

\begin{proof}
Suppose the translation orbit of $u$ is $A$-stable.  Then
\[
 a\cdot u=\tau_z u
\]
for some $z\in V$, determined modulo $V_u$.  The subgroup $V_u$ is
$T$-invariant.  Indeed, if $t\in V_u$, then
\[
 a\cdot(\tau_tu)=\tau_{Tt}(a\cdot u)=\tau_{Tt+z}u.
\]
Since $a\cdot u=\tau_z u$, it follows that $\tau_{Tt}u=u$, and hence
$Tt\in V_u$.

The operator $T$ has no nonzero fixed point on $V/V_u$.  If
$t+V_u$ is fixed, then $T^it\equiv t\pmod{V_u}$ for every $i$.  Put
\[
 s=\sum_{i=0}^{20}T^it.
\]
The vector $s$ is fixed by $T$, so $s=0$ by
\eqref{eq:no-fixed-V}.  On the other hand,
$s\equiv21t\equiv t\pmod{V_u}$, because $V$ has characteristic $2$.
Thus $t\in V_u$.

Therefore $T-I$ has trivial kernel on the finite-dimensional space
$V/V_u$, and is invertible.  Choose $w$ with
\[
 (T-I)w\equiv z\pmod{V_u}.
\]
Since $V$ has characteristic $2$, the congruence gives
$(T-I)w+z\in V_u$, and hence $Tw+z\equiv w\pmod{V_u}$.  Thus
\[
 a\cdot(\tau_wu)=\tau_{Tw+z}u=\tau_wu,
\]
so the orbit contains an $A$-fixed element.

If both $u$ and $\tau_tu$ are fixed, then $t+V_u$ is fixed in
$V/V_u$.  The preceding argument gives $t\in V_u$, so
$\tau_tu=u$.  The fixed representative is unique.
\end{proof}

\begin{proposition}\label{prop:invariant-irr}
The $A$-invariant irreducible characters of $G$ are precisely
\[
 \chi_u=\Ind_{B\rtimes V_u}^G\widetilde\lambda_u,
 \qquad u\in B^A,
\]
where
\[
 \widetilde\lambda_u(f,t)=\lambda_u(f).
\]
In particular,
\[
 |\Irr^A(G)|=4096.
\]
\end{proposition}

\begin{proof}
By Proposition~\ref{prop:all-irr}, an invariant irreducible character
must arise from an $A$-stable translation orbit.  Lemma
\ref{lem:fixed-representative} supplies a unique representative
$u\in B^A$.

For fixed $u\in B^A$, the extension $\widetilde\lambda_u$ is
$A$-invariant, and
\[
 \chi_{u,\mu}^{\,a}=\chi_{u,\mu^a}.
\]
The parameterization in Proposition~\ref{prop:all-irr} is injective, so
$\chi_{u,\mu}$ is invariant if and only if $\mu$ is invariant.

Write $\mu(t)=(-1)^{\ell(t)}$ for an $\F_2$-linear functional
$\ell$ on $V_u$.  Invariance gives $\ell(Tt)=\ell(t)$, so $\ell$
vanishes on $(T-I)V_u$.  By \eqref{eq:no-fixed-V}, the restriction of
$T-I$ to $V_u$ is injective.  Since $V_u$ is finite-dimensional, this
restriction is invertible.  Hence $\ell=0$ and $\mu$ is trivial.
The parameterization follows, and Proposition
\ref{prop:fixed-subgroup} gives $|B^A|=4096$.
\end{proof}

For $u,c\in B^A$, let $x_t=(0,t)$.  Since $I_u$ is normal in
$G$ and $c\in B\subseteq I_u$, the induced-character formula, with
$x_t$ ranging over representatives for $G/I_u$, gives
\[
 \chi_u(c)=\sum_{t\in V/V_u}
 \widetilde\lambda_u\bigl(x_t^{-1}(c,0)x_t\bigr).
\]
Now
\[
 x_t^{-1}(c,0)x_t=(\tau_{-t}c,0)
\]
and
\[
 \langle u,\tau_{-t}c\rangle=\langle\tau_tu,c\rangle.
\]
Therefore
\begin{equation}\label{eq:value-formula}
 \chi_u(c)=\sum_{t\in V/V_u}(-1)^{\langle\tau_tu,c\rangle}.
\end{equation}
The summand is constant on cosets of $V_u$, so
\begin{equation}\label{eq:value-formula-full}
 |V_u|\chi_u(c)=
 \sum_{t\in V}(-1)^{\langle\tau_tu,c\rangle}.
\end{equation}

\section{Radial correlations}

The action of $\langle\omega\rangle$ on $U$ has six orbits,
\[
 \{0\},L_1^\times,\ldots,L_5^\times.
\]
A Boolean function on $U$ that is constant on these orbits will be
called radial.  We encode it by
\[
 X=(x_0,x_1,\ldots,x_5)=(x_0,x)\in\F_2^6
\]
and write
\[
 \tau(X)=x_0+\sum_{i=1}^5x_i.
\]

\begin{lemma}[Two-slice decomposition]\label{lem:two-slice}
Every $u\in B^A$ has a unique representation
\[
 u(z,w)=
 \begin{cases}
 X(z),&w=0,\\
 Y(z),&w\neq0,
 \end{cases}
\]
where $X$ and $Y$ are radial functions on $U$.
\end{lemma}

\begin{proof}
The slice at $w=0$ is radial by invariance under $a$.  If $w\neq0$,
then $a^7$ fixes $w$ and acts on $U$ as multiplication by
$\omega^7=\omega$, so the slice at $w$ is also radial.

For $w,w'\neq0$, choose $k$ such that $\zeta^kw=w'$.  Invariance gives
\[
 u(z,w)=u(\omega^kz,w').
\]
The slice at $w'$ is radial, so the right side equals $u(z,w')$.
Hence all nonzero slices coincide.  Uniqueness is immediate.
\end{proof}

We therefore encode $u\in B^A$ by a pair
$(X,Y)\in\F_2^6\times\F_2^6$, and use $(E,F)$ for $c\in B^A$.

For radial functions $X,E$ on $U$, define
\[
 K_X(E)(s)=\sum_{z\in U}X(z+s)E(z)\in\F_2.
\]
This is again radial.

\begin{lemma}[Correlation formula]\label{lem:correlation}
Let $X=(x_0,x)$ and $E=(e_0,e)$, where $x,e\in\F_2^5$.  Put
$s_x=\sum_i x_i$, $s_e=\sum_i e_i$, and
$\one=(1,1,1,1,1)$.  Then
\begin{equation}\label{eq:correlation}
 K_X(E)=
 \left(
 x_0e_0+x\cdot e,
 \ \tau(E)x+\tau(X)e+(s_xs_e+x\cdot e)\one
 \right).
\end{equation}
\end{lemma}

\begin{proof}
At $s=0$,
\[
 K_X(E)(0)=x_0e_0+\sum_{i=1}^53x_ie_i=x_0e_0+x\cdot e.
\]
Fix $s\in L_i^\times$.  The terms $z=0$ and $z=s$ contribute
$x_ie_0+x_0e_i$.  The other two nonzero elements of $L_i$ contribute
$x_ie_i$ twice and cancel.

For $j\neq i$, the three points of $s+L_j^\times$ lie one each on the
three projective lines different from $L_i$ and $L_j$.  Indeed, none
lies on $L_i$ or $L_j$, and two cannot lie on the same third line,
since their difference would belong both to that line and to $L_j$.
The contribution of $L_j^\times$ is therefore
\[
 e_j\sum_{k\neq i,j}x_k.
\]
Thus
\[
 K_X(E)_i=x_ie_0+x_0e_i+
 \sum_{j\neq i}e_j\sum_{k\neq i,j}x_k.
\]
Expanding this expression in $\F_2$ yields the $i$th coordinate in
\eqref{eq:correlation}.
\end{proof}

Let $u=(X,Y)$ and $c=(E,F)$.  For $t=(s,v)\in U\oplus W$, put
\[
 h_{u,c}(t)=\langle\tau_tu,c\rangle.
\]
If $v=0$, then
\begin{equation}\label{eq:H0}
 H_0(s):=h_{u,c}(s,0)=K_X(E)(s)+K_Y(F)(s).
\end{equation}
If $v\neq0$, then
\begin{equation}\label{eq:H1}
 H_1(s):=h_{u,c}(s,v)=K_Y(E)(s)+K_X(F)(s).
\end{equation}
For \eqref{eq:H1}, the $W$-coordinates $0$ and $v$ interchange the two
slices.  The other six nonzero coordinates contribute the same term six
times and cancel in $\F_2$.  In particular, $H_1$ is independent of
the chosen nonzero $v$.

Set
\[
 P=X+Y,\qquad Q=E+F.
\]
Bilinearity gives
\begin{align}
 H_0+H_1&=K_P(Q),\label{eq:Hsum}\\
 H_0&=K_P(E)+K_Y(Q).\label{eq:H0-reduced}
\end{align}

For a radial function $H=(h_0,h_1,\ldots,h_5)$, define
\[
 S(H)=(-1)^{h_0}+3\sum_{i=1}^5(-1)^{h_i}.
\]
Equation \eqref{eq:value-formula-full} becomes
\begin{equation}\label{eq:value-radial}
 |V_u|\chi_u(c)=S(H_0)+7S(H_1).
\end{equation}

\section{The exact zero criterion}

Let
\[
 \delta=(1,0,0,0,0,0),\qquad
 J=(1,1,1,1,1,1).
\]
For $R\subseteq\{1,\ldots,5\}$, let
\[
 \rho_R=(0,\one_R),
\]
where $\one_R$ is the indicator vector of $R$.

\begin{lemma}[Zero patterns]\label{lem:zero-patterns}
Equation \eqref{eq:value-radial} vanishes if and only if
\begin{equation}\label{eq:zero-pattern}
 (H_0,H_1)=(\delta+\varepsilon J,\rho_R+\varepsilon J)
\end{equation}
for some $\varepsilon\in\F_2$ and some three-element subset
$R\subseteq\{1,\ldots,5\}$.
\end{lemma}

\begin{proof}
Write the signs associated with $H_0$ as
$a,b_1,\ldots,b_5\in\{\pm1\}$ and those associated with $H_1$ as
$d,e_1,\ldots,e_5\in\{\pm1\}$.  The vanishing equation is
\[
 a+3\sum_{i=1}^5b_i+7d+21\sum_{i=1}^5e_i=0.
\]
The last sum is odd.  The absolute value of the preceding terms is at
most $1+15+7=23$, so $\sum e_i=\pm1$.

If $\sum e_i=-1$, the remaining terms must total $21$.  This forces
$a=-1$, $b_i=1$ for all $i$, and $d=1$.  Exactly three of the $e_i$
are negative.  This is the pattern $(\delta,\rho_R)$ with $|R|=3$.
The case $\sum e_i=1$ is its simultaneous sign reversal, obtained by
adding $J$ to both binary vectors.
\end{proof}

A zero therefore requires
\begin{equation}\label{eq:target}
 K_P(Q)=H_0+H_1=\delta+\rho_R
\end{equation}
for a three-set $R$.

Write
\[
 P=(p_0,p),\qquad p\in\F_2^5.
\]

\begin{lemma}\label{lem:odd-augmentation}
If $\tau(P)=1$, then $K_P$ is invertible on the six-dimensional space
of radial functions on $U$.
\end{lemma}

\begin{proof}
Identify all functions on the additive group $U\cong C_2^4$ with the
group algebra $\F_2[U]$.  Since every element of $U$ is its own
inverse, the coefficient of $s$ in the product of $P$ and $Q$ is
$K_P(Q)(s)$.  Thus $K_P$ is multiplication by $P$.

The augmentation ideal of $\F_2[U]$ is nilpotent.  For example, after
choosing a basis $g_1,\ldots,g_4$ of $U$ and setting $z_i=g_i+1$, one
has
\[
 \F_2[U]\cong
 \F_2[z_1,z_2,z_3,z_4]/(z_1^2,z_2^2,z_3^2,z_4^2).
\]
An element of augmentation $1$ is therefore $1+n$ with $n$ nilpotent,
and is a unit.  The augmentation of the radial function $P$ is
\[
 p_0+3\sum_{i=1}^5p_i=\tau(P).
\]
Hence $P$ is a unit when $\tau(P)=1$.

Multiplication by $\omega$ induces an automorphism of the group algebra
that fixes $P$.  It also fixes the unique inverse of $P$.  The inverse
is therefore radial, so multiplication by $P$ restricts to an
invertible map on the radial subspace.
\end{proof}

\begin{proposition}[Complete zero criterion]\label{prop:zero-criterion}
Let $u=(X,Y)\in B^A$, put $P=X+Y=(p_0,p)$, and write
$Y=(y_0,y)$.  If $\tau(P)=0$, set
\[
 r=p+p_0\one.
\]
Then $\chi_u$ vanishes at some element of $C$ if and only if at least
one of the following holds:
\begin{enumerate}[label=\textup{(\roman*)}]
 \item $\tau(P)=1$;
 \item $\tau(P)=0$, $\wt(r)=2$, $\tau(Y)=1$, and
 \[
 r\cdot y=1+p_0.
 \]
\end{enumerate}
\end{proposition}

\begin{proof}
Suppose first that $\tau(P)=1$.  Lemma
\ref{lem:odd-augmentation} allows us to choose, for any three-set $R$,
\[
 Q=K_P^{-1}(\delta+\rho_R)
\]
and then
\[
 E=K_P^{-1}(\delta+K_Y(Q)).
\]
Equations \eqref{eq:Hsum} and \eqref{eq:H0-reduced} give
$H_0=\delta$ and $H_1=\rho_R$.  Lemma
\ref{lem:zero-patterns} yields a zero.

Assume now that $\tau(P)=0$.  For $Q=(q_0,q)$, formula
\eqref{eq:correlation} simplifies to
\begin{equation}\label{eq:image-even}
 K_P(Q)=\alpha J+\beta(0,r),
\end{equation}
where
\[
 \alpha=p_0q_0+p\cdot q,
 \qquad
 \beta=\tau(Q).
\]
Indeed, the first coordinate is $\alpha$, while the five line
coordinates are $\alpha\one+\beta r$.  The parameters satisfy
\begin{equation}\label{eq:alpha-beta}
 \alpha+p_0\beta=r\cdot q.
\end{equation}
If $r\neq0$, any prescribed pair $(\alpha,\beta)$ occurs: choose $q$
with $r\cdot q=\alpha+p_0\beta$ and then put
$q_0=\beta+\sum_iq_i$.

The target in \eqref{eq:target} has first coordinate $1$ and a line
part of weight $3$.  A vector in \eqref{eq:image-even} with first
coordinate $1$ has line part $\one$ or $\one+r$, of weights $5$ and
$5-\wt(r)$.  Since $r$ has even weight, the target occurs precisely
when $\wt(r)=2$.  In that case $R$ is the complement of $\supp(r)$,
and \eqref{eq:target} is equivalent to
\begin{equation}\label{eq:Q-critical}
 \tau(Q)=1,
 \qquad
 r\cdot q=1+p_0.
\end{equation}

We next determine when $E$ can be chosen in
\eqref{eq:H0-reduced}.  Because $J$ belongs to the image of $K_P$,
the parameter $\varepsilon$ in \eqref{eq:zero-pattern} is immaterial.
The required solvability is equivalent to
\begin{equation}\label{eq:residual-membership}
 \delta+K_Y(Q)\in\langle J,(0,r)\rangle.
\end{equation}
Write $t=\tau(Y)$ and $\sigma=\sum_iq_i$.  Since $\tau(Q)=1$, one has
$q_0=1+\sigma$.  Substituting \eqref{eq:correlation} into
\eqref{eq:residual-membership}, and subtracting the multiple of $J$
determined by the first coordinate, gives
\begin{equation}\label{eq:critical-affine}
 y+tq+(1+y_0+t\sigma)\one\in\{0,r\}.
\end{equation}

If $t=0$, the coordinate sum of the left side of
\eqref{eq:critical-affine} is
\[
 \sum_i y_i+1+y_0=1,
\]
since $\sum_i y_i=y_0$.  Both $0$ and $r$ have even coordinate sum,
so $t=0$ is impossible.  Thus $\tau(Y)=1$.

Taking the dot product of \eqref{eq:critical-affine} with $r$ gives
$r\cdot y=r\cdot q$, because $r\cdot\one=r\cdot r=0$ in $\F_2$.
Equation \eqref{eq:Q-critical} now gives
\[
 r\cdot y=1+p_0.
\]
Conversely, suppose $\tau(Y)=1$ and $r\cdot y=1+p_0$.  Define
\[
 q=y+(1+y_0)\one,
 \qquad q_0=1.
\]
Since $\sum_i y_i=1+y_0$, one has $\sum_iq_i=0$ and hence
$\tau(Q)=1$.  Also $r\cdot q=r\cdot y=1+p_0$, because
$r\cdot\one=0$.  Equation \eqref{eq:Q-critical} holds, and the left
side of \eqref{eq:critical-affine} is zero.  Thus an appropriate $E$
exists, and the character has a zero.
\end{proof}

\section{Counting the zero-free characters}

\begin{theorem}\label{thm:main}
For the action constructed in Section~2,
\[
 |N_A(G)|=1728
 \qquad\text{and}\qquad
 |C_G(A)/C_G(A)'|=4096.
\]
In particular, the equality in \eqref{eq:question} does not hold for
all coprime actions.
\end{theorem}

\begin{proof}
The second equality is Proposition~\ref{prop:fixed-subgroup}.  To
obtain the first, count the pairs $(X,Y)$ that fail both conditions in
Proposition~\ref{prop:zero-criterion}.

The change of variables
\[
 (X,Y)\longleftrightarrow(P,Y),
 \qquad P=X+Y,
\]
is a bijection on $\F_2^6\times\F_2^6$.  There are $32$ choices of
$P$ with $\tau(P)=1$, and Proposition
\ref{prop:zero-criterion} shows that all $64$ choices of $Y$ then give
vanishing characters.

Suppose $\tau(P)=0$.  The map
\[
 P=(p_0,p)\longmapsto(p_0,r),
 \qquad r=p+p_0\one,
\]
is a bijection between such $P$ and pairs in which $p_0\in\F_2$ and
$r\in\F_2^5$ has even weight.  There are
\[
 2,\qquad 2\binom52=20,\qquad 2\binom54=10
\]
choices with $\wt(r)=0,2,4$, respectively.

If $\wt(r)=0$ or $4$, every one of the $64$ choices of $Y$ is
zero-free.  Fix a critical $P$ with $\wt(r)=2$.  Vanishing is equivalent
to the two affine equations
\[
 \tau(Y)=1,
 \qquad
 r\cdot y=1+p_0.
\]
These equations are independent: the first involves $y_0$, while the
second does not and $r\neq0$.  Exactly $2^{6-2}=16$ of the $64$ choices
of $Y$ satisfy both equations, leaving $48$ zero-free choices.
Therefore
\[
 \begin{aligned}
 |N_A(G)|
 &=2\cdot64+10\cdot64+20\cdot48\\
 &=128+640+960\\
 &=1728.
 \end{aligned}
\]
Since $1728<4096$, the equality in \eqref{eq:question} fails.
\end{proof}

\begin{remark}[A single explicit zero]\label{rem:witness}
The strict inequality can be seen without the full count.  Let
$u=\delta_0$ be the indicator of $0\in V$.  Then $V_u=0$, and
$\chi_u=\Ind_B^G\lambda_u$ is an $A$-invariant irreducible character of
degree $128$.  Choose three distinct projective lines
$L_1,L_2,L_3$ in $U$ and let $c$ be the indicator of
\[
 \{0\}\cup\bigcup_{i=1}^3(L_i^\times\times W^\times).
\]
This set is $A$-stable and has size $1+3\cdot21=64$, so $c\in C$ and
\[
 \chi_u(c)=\sum_{t\in V}(-1)^{c(t)}=64-64=0.
\]
The full analysis above gives the exact count $1728$.
\end{remark}

\section{Verification and limitations}

Three exact Python implementations were used as independent checks: a
full integer Walsh-transform computation, a direct computation in
$\F_4^2\oplus\F_8$, and a binary linear-algebra test based on the twenty
zero patterns in Lemma~\ref{lem:zero-patterns}.  Each returns $1728$
zero-free characters and $2368$ characters with a zero.  These
computations are not used in the proof.  The source code is available at
\url{https://github.com/erichou1/coprime-action-character-count-counterexample}.

The example is not claimed to be minimal.  It uses a cyclic operator
group of order $21$, and the argument does not address versions of the
question with additional restrictions on $A$ or $G$.

\section*{Acknowledgments}

The author used OpenAI's ChatGPT (GPT-5.6 Thinking) during the
preparation of this manuscript for mathematical exploration, checking
intermediate arguments and edge cases, generating exact verification
code, and editorial revision.  The author reviewed the resulting
arguments and takes responsibility for the contents of the paper.

\begin{thebibliography}{9}

\bibitem{kourovka}
E. I. Khukhro and V. D. Mazurov (eds.),
\emph{Unsolved Problems in Group Theory: The Kourovka Notebook},
No.~21, Sobolev Institute of Mathematics, Novosibirsk, 2026,
Problem~21.100, p.~178; arXiv:1401.0300v45.

\end{thebibliography}

\end{document}
